\documentclass[../libro.tex]{subfiles}

\begin{document}

\ifSubfilesClassLoaded{\mainmatter\chapter{行列式}\clearpage}{}

\section{阵的积}

前面, 我们学习了阵的运算, 如阵的加、减、数乘.
本节, 我们学习阵的一个新运算, 阵的乘法.
对若干个阵作乘法的结果就是这些阵的积.

\begin{definition}
    设 \(A\) 是 \(m \times s\)~阵,
    且 \(B\) 是 \(s \times n\)~阵
    (也就是, \(A\) 有多少列, \(B\) 就有多少行).
    我们定义 \(A\) 与 \(B\)~的\emph{积}
    (注意 \(A\), \(B\) 的次序) 为%
    一个 \(m \times n\)~阵 \(AB\),
    其中, 对任何不超过 \(m\) 的正整数 \(i\)
    与任何不超过 \(n\) 的正整数 \(j\),
    \begin{align*}
        [AB]_{i,j}
        = {} &
        [A]_{i,1} [B]_{1,j}
        + [A]_{i,2} [B]_{2,j}
        + \dots
        + [A]_{i,s} [B]_{s,j}
        \\
        = {} &
        \sum_{k = 1}^{s} {[A]_{i,k} [B]_{k,j}}.
    \end{align*}
    (通俗地, \(AB\) 是 \(m \times n\)~阵,
    其的 \((i, j)\)-元是
    \(A\)~的行~\(i\) 跟 \(B\)~的列~\(j\)~的%
    相应位置的元的积的和.)
\end{definition}

可以看到, 阵的乘法跟阵的加法比, 有大的区别.
首先, 加法要求二个阵的尺寸相同, 而乘法要求%
第~1~个阵的列数等于第~2~个阵的行数.
其次, 阵的加法的定义用到的只不过是数的加法,
而阵的乘法的定义要同时用到数的加法与乘法.

阵的乘法, 不像数的乘法, 有可能是不可换的.
具体地,
设 \(A\), \(B\) 分别是 \(m \times s\) 与 \(s \times n\)~阵.
根据定义, \(AB\) 是有意义的.
不过, \(BA\) 无意义, 除非 \(n = m\).
这还只是一方面.
另一方面, 就算 \(n = m\),
\(AB\) 跟 \(BA\)~的尺寸也不一样,
除非 \(m = s = n\).
最后, 就算 \(m = s = n\),
\(AB\) 也可以不等于 \(BA\).

\begin{example}
    设
    \begin{align*}
        A = \begin{bmatrix}
                3  & 2 \\
                10 & 7 \\
            \end{bmatrix},
        \quad
        B = \begin{bmatrix}
                5  & 4 \\
                11 & 9 \\
            \end{bmatrix}.
    \end{align*}
    根据定义,
    \begin{align*}
         &
        AB = \begin{bmatrix}
                 3 \cdot 5 + 2 \cdot 11  & 3 \cdot 4 + 2 \cdot 9  \\
                 10 \cdot 5 + 7 \cdot 11 & 10 \cdot 4 + 7 \cdot 9 \\
             \end{bmatrix}
        = \begin{bmatrix}
              37  & 30  \\
              127 & 103 \\
          \end{bmatrix},
        \\
         &
        BA = \begin{bmatrix}
                 5 \cdot 3 + 4 \cdot 10  & 5 \cdot 2 + 4 \cdot 7  \\
                 11 \cdot 3 + 9 \cdot 10 & 11 \cdot 2 + 9 \cdot 7 \\
             \end{bmatrix}
        = \begin{bmatrix}
              55  & 38 \\
              123 & 85 \\
          \end{bmatrix}.
    \end{align*}
    所以, 即使 \(AB\), \(BA\) 都是 \(2\)~级阵,
    它们也不相等.
\end{example}

我们知道, 若 \(a\), \(b\) 是二个非零的数,
那么 \(ab \neq 0\).
不过, 对二个阵 \(A\), \(B\),
即使 \(A \neq 0\), 且 \(B \neq 0\),
仍有可能 \(AB = 0\)
(其中, \(0\) 是零阵,
即元全为 \(0\) 的阵).

\begin{example}
    设
    \begin{align*}
        A = \begin{bmatrix}
                0 & 0 \\
                1 & 0 \\
            \end{bmatrix},
        \quad
        B = \begin{bmatrix}
                0 & 0 \\
                0 & 1 \\
            \end{bmatrix}.
    \end{align*}
    根据定义,
    \begin{align*}
         &
        AB = \begin{bmatrix}
                 0 \cdot 0 + 0 \cdot 0 & 0 \cdot 0 + 0 \cdot 1 \\
                 1 \cdot 0 + 0 \cdot 0 & 1 \cdot 0 + 0 \cdot 1 \\
             \end{bmatrix}
        = \begin{bmatrix}
              0 & 0 \\
              0 & 0 \\
          \end{bmatrix}.
    \end{align*}
\end{example}

阵的乘法跟数的乘法虽有不同的地方,
但更重要地, 有相同的 (或类似的) 地方.

\begin{theorem}
    设 \(A\), \(B\), \(C\) 分别是
    \(m \times s\), \(s \times n\), \(n \times p\)~阵.
    设 \(D\) 是 \(m \times s\)~阵.
    设 \(G\) 是 \(s \times n\)~阵.
    设 \(x\) 是数.
    则:

    (1)
    (结合律)
    \((AB)C = A(BC)\);

    (2)
    (分配律~1)
    \((A + D) B = AB + DB\);

    (3)
    (分配律~2)
    \(A (B + G) = AB + AG\);

    (4)
    (单位阵的作用)
    \(I_m A = A = A I_s\),
    其中, \(I_m\), \(I_s\) 分别是 \(m\), \(s\)~级单位阵;

    (5)
    (阵的数乘与乘法)
    \((xA)B = x(AB) = A(xB)\).
\end{theorem}

\begin{proof}
    (1)
    注意, \(AB\) 是 \(m \times n\)~阵,
    故 \((AB)C\) 是 \(m \times p\)~阵;
    注意, \(BC\) 是 \(s \times p\)~阵,
    故 \(A(BC)\) 也是 \(m \times p\)~阵.
    根据阵的积的定义,
    \begin{align*}
        [(AB)C]_{i,j}
        = {} &
        \sum_{\ell = 1}^{n} {[AB]_{i,\ell} [C]_{\ell,j}}
        \\
        = {} &
        \sum_{\ell = 1}^{n} {
        \left(\sum_{k = 1}^{s} {[A]_{i,k} [B]_{k,\ell}} \right)
        [C]_{\ell,j} }
        \\
        = {} &
        \sum_{\ell = 1}^{n} {
        \sum_{k = 1}^{s} {
        [A]_{i,k} [B]_{k,\ell} [C]_{\ell,j} }}
        \\
        = {} &
        \sum_{k = 1}^{s} {
        \sum_{\ell = 1}^{n} {
        [A]_{i,k} [B]_{k,\ell} [C]_{\ell,j}}}
        \\
        = {} &
        \sum_{k = 1}^{s} {
        [A]_{i,k}
        \left(
        \sum_{\ell = 1}^{n} {[B]_{k,\ell} [C]_{\ell,j}}
        \right) }
        \\
        = {} &
        \sum_{k = 1}^{s} {[A]_{i,k} [BC]_{k,j}}
        \\
        = {} &
        [A(BC)]_{i,j}.
    \end{align*}

    (2)
    您还是要先说明, 等式二侧的阵的尺寸都是 \(m \times n\)
    (我留它为您的习题).
    根据阵的和与积的定义,
    \begin{align*}
        [(A + D)B]_{i,j}
        = {} &
        \sum_{k = 1}^{s} {[A + D]_{i,k} [B]_{k,j}}
        \\
        = {} &
        \sum_{k = 1}^{s} {([A]_{i,k} + [D]_{i,k})\, [B]_{k,j}}
        \\
        = {} &
        \sum_{k = 1}^{s} {([A]_{i,k} [B]_{k,j}
        + [D]_{i,k} [B]_{k,j})}
        \\
        = {} &
        \sum_{k = 1}^{s} {[A]_{i,k} [B]_{k,j}}
        +
        \sum_{k = 1}^{s} {[D]_{i,k} [B]_{k,j}}
        \\
        = {} &
        [AB]_{i,j} + [DB]_{i,j}
        \\
        = {} &
        [AB + DB]_{i,j}.
    \end{align*}

    (3)
    这跟分配律~1 的论证完全类似;
    我留它为您的习题.

    (4)
    先证 \(I_m A = A\).
    您还是要先说明, 等式二侧的阵的尺寸都是 \(m \times s\)
    (我留它为您的习题).
    回想, 单位阵 \(I\) 的 \((i, j)\)-元是
    \(\delta(i, j)\);
    具体地, \(i = j\) 时, 它是 \(1\),
    而 \(i \neq j\) 时, 它是 \(0\).
    根据阵的积的定义,
    \begin{align*}
        [I_m A]_{i,j}
        = {} &
        \sum_{u = 1}^{m} {[I_m]_{i,u} [A]_{u,j}}
        \\
        = {} &
        \sum_{u = 1}^{m} {\delta(i, u)\, [A]_{u,j}}
        \\
        = {} &
        \delta(i, i)\, [A]_{i,j}
        + \sum_{\substack{1 \leq u \leq m \\ u \neq i}}
        {\delta(i, u)\, [A]_{u,j}}
        \\
        = {} &
        1\,[A]_{i,j}
        + \sum_{\substack{1 \leq u \leq m \\ u \neq i}}
        {0\,[A]_{u,j}}
        \\
        = {} & [A]_{i,j}.
    \end{align*}
    请允许我留另一部分为您的习题.

    (5)
    先证 \((xA)B = x(AB)\).
    您还是要先说明, 等式二侧的阵的尺寸都是 \(m \times n\)
    (我留它为您的习题).
    根据阵的数乘与阵的积的定义,
    \begin{align*}
        [(xA)B]_{i,j}
        = {} &
        \sum_{k = 1}^{s} {[xA]_{i,k} [B]_{k,j}}
        \\
        = {} &
        \sum_{k = 1}^{s} {x[A]_{i,k} [B]_{k,j}}
        \\
        = {} &
        x \sum_{k = 1}^{s} {[A]_{i,k} [B]_{k,j}}
        \\
        = {} &
        x [AB]_{i,j}
        \\
        = {} &
        [x(AB)]_{i,j}.
    \end{align*}
    请允许我留另一部分为您的习题.
\end{proof}

由此, 我们可简单地写 \((AB)C\) 或 \(A(BC)\) 为 \(ABC\),
且可简单地写 \(x(AB)\), \((xA)B\), \(A(xB)\)
为 \(xAB\)
(当然, \(x\) 是一个数).

\begin{theorem}
    设 \(B\) 是 \(s \times m\)~阵.
    设 \(A\) 是 \(m \times n\)~阵.
    设 \(A\) 的列~\(1\), \(2\), \(\dots\), \(n\)
    分别是 \(a_1\), \(a_2\), \(\dots\), \(a_n\)
    (于是, 我们可写 \(A = [a_1, a_2, \dots, a_n]\)).
    则
    \begin{align*}
        BA = [Ba_1, Ba_2, \dots, Ba_n].
    \end{align*}
\end{theorem}

\begin{proof}
    每个 \(Ba_j\) 都是 \(s \times 1\)~阵,
    故 \([Ba_1, Ba_2, \dots, Ba_n]\) 是 \(s \times n\)~阵.
    当然, \(BA\) 也是 \(s \times n\)~阵.
    根据阵的积的定义,
    \begin{align*}
        [BA]_{i,j}
        = {} & \sum_{k = 1}^{m} {[B]_{i,k} [A]_{k,j}}
        \\
        = {} & \sum_{k = 1}^{m} {[B]_{i,k} [a_j]_{k,1}}
        \\
        = {} & [Ba_j]_{i,1}
        \\
        = {} & [\, [Ba_1, Ba_2, \dots, Ba_n] \,]_{i,j}.
        \qedhere
    \end{align*}
\end{proof}

\begin{theorem}
    设 \(A\) 是 \(m \times n\)~阵.
    设 \(A\) 的列~\(1\), \(2\), \(\dots\), \(n\)
    分别是 \(a_1\), \(a_2\), \(\dots\), \(a_n\)
    (于是, 我们可写 \(A = [a_1, a_2, \dots, a_n]\)).
    设 \(X\) 是一个 \(n \times 1\)~阵.
    则
    \begin{align*}
        AX = [X]_{1,1} a_1 + [X]_{2,1} a_2 + \dots
        + [X]_{n,1} a_n.
    \end{align*}
\end{theorem}

\begin{proof}
    \(AX\) 自然是一个 \(m \times 1\)~阵.
    不难看出, 待证等式的右侧也是 \(m \times 1\)~阵.
    根据阵的积、和、数乘的定义,
    \begin{align*}
        [AX]_{i,1}
        = {} & \sum_{k = 1}^{n} {[A]_{i,k} [X]_{k,1}}
        \\
        = {} & \sum_{k = 1}^{n} {[a_k]_{i,1} [X]_{k,1}}
        \\
        = {} & \sum_{k = 1}^{n} {[[X]_{k,1} a_k]_{i,1}}
        \\
        = {} & \left[ \sum_{k = 1}^{n} {[X]_{k,1}
        a_k } \right]_{i,1}.
        \qedhere
    \end{align*}
\end{proof}

\section{Binet--Cauchy 公式 (青春版)}

设 \(A\), \(B\) 都是 \(n\)~级阵.
那么, \(BA\) 当然也是 \(n\)~级阵.
\(A\), \(B\), \(BA\) 都是 \(n\)~级阵,
故, 它们都有行列式.
本节, 我们学习三者的行列式的关系.

\begin{theorem}[Binet--Cauchy \pinjino{bǐnèi--kōxi} 公式, 青春版]
    设 \(A\), \(B\) 都是 \(n\)~级阵.
    则
    \begin{align*}
        \det {(BA)} = \det {(B)} \det {(A)}.
    \end{align*}
\end{theorem}

我用一个例助您理解, 此定理在说什么.

\begin{example}
    设
    \begin{align*}
        A = \begin{bmatrix}
                3  & 2 \\
                10 & 7 \\
            \end{bmatrix},
        \quad
        B = \begin{bmatrix}
                5  & 4 \\
                11 & 9 \\
            \end{bmatrix}.
    \end{align*}
    不难算出
    \begin{align*}
        AB
        = \begin{bmatrix}
              37  & 30  \\
              127 & 103 \\
          \end{bmatrix},
        \quad
        BA
        = \begin{bmatrix}
              55  & 38 \\
              123 & 85 \\
          \end{bmatrix}.
    \end{align*}
    可以看到, \(AB \neq BA\).

    不过, \(\det {(AB)} = \det {(BA)}\).
    一方面, 我们可直接验证:
    \begin{align*}
         & \det {(AB)} = 37 \cdot 103 - 127 \cdot 30 = 1; \\
         & \det {(BA)} = 55 \cdot 85 - 123 \cdot 38 = 1.
    \end{align*}
    另一方面, Binet--Cauchy 公式 (青春版) 指出,
    \begin{align*}
        \det {(AB)}
        = {} & \det {(A)} \det {(B)} \\
        = {} & \det {(B)} \det {(A)} \\
        = {} & \det {(BA)},
    \end{align*}
    因为数的乘法是可换的.
\end{example}

\begin{proof}
    考虑定义在全体 \(n\)~级阵上的函数
    \(f(C) = \det {(BC)}\),
    其中, \(C = [c_1, c_2, \dots, c_n]\) 是 \(n\)~级阵.

    (1)
    \(f\) 是多线性的.
    对任何不超过 \(n\) 的正整数 \(j\),
    任何 \(n-1\)~个 \(n \times 1\)~阵
    \(c_1\), \(\dots\), \(c_{j-1}\),
    \(c_{j+1}\), \(\dots\), \(c_n\),
    任何二个 \(n \times 1\)~阵 \(x\), \(y\),
    任何二个数 \(s\), \(t\),
    有
    \begin{align*}
             &
        f([\dots, c_{j-1}, sx + ty, c_{j+1}, \dots])
        \\
        = {} &
        \det {[\dots, Bc_{j-1}, B(sx + ty), Bc_{j+1}, \dots]}
        \\
        = {} &
        \det {[\dots, Bc_{j-1}, sBx + tBy, Bc_{j+1}, \dots]}
        \\
        = {} &
        s\det {[\dots, Bc_{j-1}, Bx, Bc_{j+1}, \dots]}
        +
        t\det {[\dots, Bc_{j-1}, By, Bc_{j+1}, \dots]}
        \\
        = {} &
        s f([\dots, c_{j-1}, x, c_{j+1}, \dots])
        +
        t f([\dots, c_{j-1}, y, c_{j+1}, \dots]).
    \end{align*}
    我们用到了阵的运算律与行列式的多线性.

    (2)
    \(f\) 是交错性的.
    因为若 \(c_1\), \(c_2\), \(\dots\), \(c_n\) 中有二个相等,
    则 \(Bc_1\), \(Bc_2\), \(\dots\), \(Bc_n\) 中也有二个相等.
    再利用行列式的交错性,
    \(f(C) = \det {(BC)} = 0\).

    所以, 对任何 \(n\)~级阵 \(A\),
    \(f(A) = f(I) \det {(A)}\).
    注意, \(BI = B\),
    故
    \begin{equation*}
        \det {(BA)} = f(A) = f(I) \det {(A)}
        = \det {(B)} \det {(A)}.
        \qedhere
    \end{equation*}
\end{proof}

或许, (方) 阵的行列式与阵的积的定义是较复杂的,
跟阵的转置、加、减、数乘比.
不过, Binet--Cauchy 公式 (的青春版) 给出了一个关系:
二个同级的方阵的积的行列式%
等于这二个阵的行列式的积.

\section{按一列 (行) 展开行列式的公式的变体}

我们回想,
我们既可按任何一列, 也可按任何一行,
展开一个阵的行列式:

% subfile fix: repeat restatable
% see:
% determinanto - determinantoj.tex
\ifSubfilesClassLoaded{%
    \begin{restatable}{theorem}{TheoremExpansionAboutAnyColumn}
        设 \(A\) 是 \(n\)~级阵 (\(n \geq 1\)).
        设 \(j\) 为整数, 且 \(1 \leq j \leq n\).
        则
        \begin{align*}
            \det {(A)} = \sum_{i = 1}^{n}
            {(-1)^{i+j} [A]_{i,j} \det {(A(i|j))}}.
        \end{align*}
    \end{restatable}
}{\TheoremExpansionAboutAnyColumn*}

% subfile fix: repeat restatable
% see:
% determinanto - ecoj de determinantoj.tex
\ifSubfilesClassLoaded{%
    \begin{restatable}{theorem}{TheoremExpansionAboutAnyRow}
        设 \(A\) 是 \(n\)~级阵 (\(n \geq 1\)).
        设 \(i\) 为整数, 且 \(1 \leq i \leq n\).
        则
        \begin{align*}
            \det {(A)} = \sum_{j = 1}^{n}
            {(-1)^{i+j} [A]_{i,j} \det {(A(i|j))}}.
        \end{align*}
    \end{restatable}
}{\TheoremExpansionAboutAnyRow*}

或许, 您还记得约定:
当 \(A\) 是 \(1\)~级阵时,
形如 \(A(1|1)\) 的记号表示 ``\(0\)~级阵'',
且 ``\(0\)~级阵'' 的行列式为 \(1\).
此约定可使我们简单地写公式.

若我们改变公式的几个文字,
则可得到不一样的结果:

\begin{theorem}
    设 \(A\) 是 \(n\)~级阵 (\(n \geq 1\)).
    设 \(k\), \(j\) 都是不超过 \(n\) 的正整数,
    且 \(k \neq j\).
    则
    \begin{align*}
        \sum_{\ell = 1}^{n}
        {(-1)^{\ell+k} [A]_{\ell,j} \det {(A(\ell|k))}}
        = 0.
    \end{align*}
    类似地, 若 \(i\), \(k\) 都是不超过 \(n\) 的正整数,
    且 \(i \neq k\),
    则
    \begin{align*}
        \sum_{s = 1}^{n}
        {(-1)^{k+s} [A]_{i,s} \det {(A(k|s))}}
        = 0.
    \end{align*}
\end{theorem}

我用一个例助您理解, 此定理在说什么.

\begin{example}
    设
    \begin{align*}
        A =
        \begin{bmatrix}
            1 & 4 & 8 \\
            2 & 6 & 5 \\
            3 & 9 & 7 \\
        \end{bmatrix}.
    \end{align*}
    不难算出
    \begin{align*}
         & \det {(A(1|1))} = 6 \cdot 7 - 9 \cdot 5 = -3;  \\
         & \det {(A(2|1))} = 4 \cdot 7 - 9 \cdot 8 = -44; \\
         & \det {(A(3|1))} = 4 \cdot 5 - 6 \cdot 8 = -28.
    \end{align*}
    那么, 这个定理说, 当 \(j = 2\) 或 \(j = 3\) 时,
    \begin{align*}
         & \hphantom{{} + {}}
        (-1)^{1+1} [A]_{1,j} \det {(A(1|1))}
        + (-1)^{2+1} [A]_{2,j} \det {(A(2|1))}
        \\
         &
        + (-1)^{3+1} [A]_{3,j} \det {(A(3|1))} = 0,
    \end{align*}
    即
    \begin{align*}
        -3\,[A]_{1,j} + 44\,[A]_{2,j} - 28\,[A]_{3,j} = 0.
    \end{align*}

    我们直接地验证此事.
    取 \(j = 2\), 有
    \begin{align*}
        (-3) \cdot 4 + 44 \cdot 6 - 28 \cdot 9 = -12 + 264 - 252 = 0.
    \end{align*}
    取 \(j = 3\), 有
    \begin{align*}
        (-3) \cdot 8 + 44 \cdot 5 - 28 \cdot 7 = -24 + 220 - 196 = 0.
    \end{align*}
\end{example}

\begin{proof}
    我证第~1~个; 我留第~2~个为您的习题.

    取 \(n\)~个数 \(z_1\), \(z_2\), \(\dots\), \(z_n\).
    我们作 \(n\)~级阵 \(B\), 其中,
    \begin{align*}
        [B]_{\ell, v} =
        \begin{cases}
            [A]_{\ell, v}, & v \neq k; \\
            z_\ell,        & v = k.
        \end{cases}
    \end{align*}
    (通俗地, \(B\)~的列~\(v\) 等于 \(A\)~的列~\(v\)
    (\(v \neq k\)),
    而 \(B\)~的列~\(k\) 的元分别是
    \(z_1\), \(z_2\), \(\dots\), \(z_n\).)
    由此可见, \(B(\ell|k) = A(\ell|k)\)
    (\(\ell = 1\), \(2\), \(\dots\), \(n\)).
    所以
    \begin{align*}
        \det {(B)}
        = {} & \sum_{\ell = 1}^{n}
        {(-1)^{\ell + k} [B]_{\ell, k} \det {(B(\ell|k))}}
        \\
        = {} & \sum_{\ell = 1}^{n}
        {(-1)^{\ell + k} z_\ell \det {(A(\ell|k))}}.
    \end{align*}
    现在, 特别地, 我们取
    \(z_1\), \(z_2\), \(\dots\), \(z_n\)
    为
    \([A]_{1,j}\), \([A]_{2,j}\), \(\dots\), \([A]_{n,j}\).
    此时, \(B\) 有相同的二列
    (\(B\)~的列~\(j\) 即为 \(B\)~的列~\(k\),
    并注意, \(k \neq j\)).
    所以 \(\det {(B)} = 0\).
    故
    \begin{equation*}
        \sum_{\ell = 1}^{n}
        {(-1)^{\ell+k} [A]_{\ell,j} \det {(A(\ell|k))}}
        = 0.
        \qedhere
    \end{equation*}
\end{proof}

用 \(\delta\)-记号, 有

\begin{theorem}
    设 \(A\) 是 \(n\)~级阵 (\(n \geq 1\)).
    若 \(k\), \(j\) 都是不超过 \(n\) 的正整数,
    则
    \begin{align*}
        \sum_{\ell = 1}^{n}
        {(-1)^{\ell+k} \det {(A(\ell|k))}\, [A]_{\ell,j}}
        = \det {(A)}\, \delta(k, j).
    \end{align*}
    类似地, 若 \(i\), \(k\) 都是不超过 \(n\) 的正整数,
    则
    \begin{align*}
        \sum_{s = 1}^{n}
        {[A]_{i,s} (-1)^{k+s} \det {(A(k|s))}}
        = \det{(A)}\, \delta(i, k).
    \end{align*}
\end{theorem}

\begin{proof}
    请允许我留它为您的习题.
    分类讨论即可
    (\(k = j\);
    \(k \neq j\);
    \(i = k\);
    \(i \neq k\)).
\end{proof}

\begin{definition}[古伴]
    设 \(A\) 是 \(n\)~级阵 (\(n \geq 1\)).
    定义 \(A\)~的\emph{古伴}
    (或者, 古典伴随阵)
    为 \(n\)~级阵 \(\operatorname{adj} {(A)}\),
    其中, 对任何不超过 \(n\) 的正整数 \(i\), \(j\),
    \begin{align*}
        [\operatorname{adj} {(A)}]_{i,j}
        = (-1)^{j+i} \det {(A(j|i))}.
    \end{align*}
    (注意等式右侧的 \(i\), \(j\) 的次序.)
\end{definition}

\(\operatorname{adj}\) 来自 \angla{adjugate} (英语).

有一件小事值得提.
根据约定,
一个 \(1\)~级阵 \(A = [a]\) 的古伴
\(\operatorname{adj} {(A)} =
[(-1)^{1+1} \det {(A(1|1))}] = [1]\),
即 \(1\)~级单位阵.

\begin{example}
    设
    \begin{align*}
        A =
        \begin{bmatrix}
            1 & 4 & 8 \\
            2 & 6 & 5 \\
            3 & 9 & 7 \\
        \end{bmatrix}.
    \end{align*}
    我们计算 \(A\)~的古伴 \(\operatorname{adj} {(A)}\).
    根据定义, 我们要计算 \(9\)~个 \(2\)~级阵的行列式:
    \begin{align*}
         & \det {(A(1|1))} = 6 \cdot 7 - 9 \cdot 5 = -3;  \\
         & \det {(A(1|2))} = 2 \cdot 7 - 3 \cdot 5 = -1;  \\
         & \det {(A(1|3))} = 2 \cdot 9 - 3 \cdot 6 = 0;   \\
         & \det {(A(2|1))} = 4 \cdot 7 - 9 \cdot 8 = -44; \\
         & \det {(A(2|2))} = 1 \cdot 7 - 3 \cdot 8 = -17; \\
         & \det {(A(2|3))} = 1 \cdot 9 - 3 \cdot 4 = -3;  \\
         & \det {(A(3|1))} = 4 \cdot 5 - 6 \cdot 8 = -28; \\
         & \det {(A(3|2))} = 1 \cdot 5 - 2 \cdot 8 = -11; \\
         & \det {(A(3|3))} = 1 \cdot 6 - 2 \cdot 4 = -2.
    \end{align*}
    所以
    \begin{align*}
             & \operatorname{adj} {(A)}
        \\
        = {} &
        \begin{bmatrix}
            (-1)^{1+1} \det {(A(1|1))}
             & (-1)^{2+1} \det {(A(2|1))}
             & (-1)^{3+1} \det {(A(3|1))} \\
            (-1)^{1+2} \det {(A(1|2))}
             & (-1)^{2+2} \det {(A(2|2))}
             & (-1)^{3+2} \det {(A(3|2))} \\
            (-1)^{1+3} \det {(A(1|3))}
             & (-1)^{2+3} \det {(A(2|3))}
             & (-1)^{3+3} \det {(A(3|3))} \\
        \end{bmatrix}
        \\
        = {} &
        \begin{bmatrix}
            +\det {(A(1|1))}
             & -\det {(A(2|1))}
             & +\det {(A(3|1))} \\
            -\det {(A(1|2))}
             & +\det {(A(2|2))}
             & -\det {(A(3|2))} \\
            +\det {(A(1|3))}
             & -\det {(A(2|3))}
             & +\det {(A(3|3))} \\
        \end{bmatrix}
        \\
        = {} &
        \begin{bmatrix}
            -3 & 44  & -28 \\
            1  & -17 & 11  \\
            0  & 3   & -2  \\
        \end{bmatrix}.
    \end{align*}
\end{example}

\begin{theorem}
    设 \(A\) 是 \(n\)~级阵 (\(n \geq 1\)).
    设 \(\operatorname{adj} {(A)}\) 为它的古伴.
    则
    \begin{align*}
        \operatorname{adj} {(A)}\, A
        = \det {(A)}\, I
        = A \operatorname{adj} {(A)}.
    \end{align*}
\end{theorem}

\begin{proof}
    首先, 二个等式的左右二侧都是 \(n\)~级阵.

    由上个定理 (的前一部分),
    \begin{align*}
        [\operatorname{adj} {(A)}\, A]_{i,j}
        = {} &
        \sum_{\ell = 1}^{n}
        {[\operatorname{adj} {(A)}]_{i,\ell} [A]_{\ell,j}}
        \\
        = {} &
        \sum_{\ell = 1}^{n}
        {(-1)^{\ell+i} \det {(A(\ell|i))}\, [A]_{\ell,j}}
        \\
        = {} & \det {(A)}\, \delta(i, j)
        \\
        = {} & \det {(A)}\, [I]_{i,j}
        \\
        = {} & [\det {(A)}\, I]_{i,j}.
    \end{align*}
    同理,
    \begin{align*}
        [A \operatorname{adj} {(A)}]_{i,j}
        = {} &
        \sum_{s = 1}^{n}
        {[A]_{i,s} [\operatorname{adj} {(A)}]_{s,j}}
        \\
        = {} &
        \sum_{s = 1}^{n}
        {[A]_{i,s} (-1)^{j+s} \det {(A(j|s))}}
        \\
        = {} & \det {(A)}\, \delta(i, j)
        \\
        = {} & \det {(A)}\, [I]_{i,j}
        \\
        = {} & [\det {(A)}\, I]_{i,j}.
        \qedhere
    \end{align*}
\end{proof}

\begin{example}
    设
    \begin{align*}
        A =
        \begin{bmatrix}
            1 & 4 & 8 \\
            2 & 6 & 5 \\
            3 & 9 & 7 \\
        \end{bmatrix}.
    \end{align*}
    我们知道, \(A\) 的行列式是 \(1\).
    我们还知道, \(A\)~的古伴
    \begin{align*}
        \operatorname{adj} {(A)} =
        \begin{bmatrix}
            -3 & 44  & -28 \\
            1  & -17 & 11  \\
            0  & 3   & -2  \\
        \end{bmatrix}.
    \end{align*}
    可以验证
    \begin{align*}
        \operatorname{adj} {(A)}\, A
        = 1I
        = A\operatorname{adj} {(A)}.
    \end{align*}
\end{example}

设 \(A\) 是 \(n\)~级阵,
且 \(\det {(A)} \neq 0\).
作 \(n\)~级阵
\(B = (\det {(A)})^{-1} \operatorname{adj} {(A)}\).
则
\begin{align*}
    BA
    = {} & ( (\det {(A)})^{-1} \operatorname{adj} {(A)} ) A
    \\
    = {} & (\det {(A)})^{-1} (\operatorname{adj} {(A)}\, A)
    \\
    = {} & (\det {(A)})^{-1} (\det {(A)}\, I)
    \\
    = {} & ( (\det {(A)})^{-1} \det {(A)} ) I
    \\
    = {} & 1I
    \\
    = {} & I.
\end{align*}
类似地, 可知 \(AB = I\).
所以, 我们有

\begin{theorem}
    设 \(A\) 是 \(n\)~级阵 (\(n \geq 1\)).
    若 \(\det {(A)} \neq 0\),
    则存在 \(n\)~级阵
    \(B = (\det {(A)})^{-1} \operatorname{adj} {(A)}\)
    使
    \begin{align*}
        BA = I = AB.
    \end{align*}
\end{theorem}

反过来, 设 \(A\) 是 \(n\)~级阵,
且有 \(n\)~级阵 \(C\) (或 \(D\))
使 \(CA = I\) (或 \(AD = I\)).
因为二个同级的方阵的积的行列式%
等于这二个阵的行列式的积,
\(1 = \det {(I)} = \det {(CA)} = \det {(C)} \det {(A)}\)
(或 \(1 = \det {(I)} = \det {(AD)} = \det {(A)} \det {(D)}\)).
从而 \(\det {(A)} \neq 0\).
则有 \(n\)~级阵
\(B = (\det {(A)})^{-1} \operatorname{adj} {(A)}\)
使 \(BA = I = AB\).
则 \(C = CI = C(AB) = (CA)B = IB = B
= (\det {(A)})^{-1} \operatorname{adj} {(A)}\)
(或 \(D = ID = (BA)D = B(AD) = BI = B
= (\det {(A)})^{-1} \operatorname{adj} {(A)}\)).

我们知道, 对任何不等于 \(0\) 的数 \(a\),
必有一个数 \(b = a^{-1}\) 使 \(ba = 1 = ab\);
反过来, 若有数 \(c\) (或 \(d\))
使 \(ca = 1\) (或 \(ad = 1\)),
则 \(a \neq 0\),
且 \(c = a^{-1}\) (或 \(d = a^{-1}\)).
这么看来,
行列式非零的阵跟非零的数是像的.

\end{document}

This material is mainly about matrix multiplication.
Matrix multiplication is a very important operation;
it will be used to discuss systems of linear equations.
The fact that
the determinant of the product of two square matrices
of the same size
is equal to the product of determinants of these two matrices
is proved,
with the theorem of determining
constant multiples of the determinant function.
The adjugate of a square matrix is also defined,
and one important result about adjugates
is immediately given.
Adjugates will be often used
in my treatment of systems of linear equations.
